\chapter{Module 1 : accès, identité et points d'entrée par rôle}

\section*{Introduction}
\addcontentsline{toc}{section}{Introduction}
Ce premier module pose le socle d'usage de toute la plateforme. Son rôle ne se limite pas à afficher un formulaire de connexion : il formalise la frontière entre les utilisateurs autorisés et les espaces métiers, sécurise l'accès aux données de recrutement et prépare l'intégration des fonctionnalités transverses comme le chat analytique. Sans cette couche, les modules de gestion des CV, du pipeline candidat et de l'IA ne pourraient pas être exposés de manière fiable dans un contexte d'entreprise.

Dans \projectname{}, l'identité n'est pas un sujet périphérique. Elle détermine les données visibles, les actions autorisées, la navigation initiale et même le périmètre de l'agent conversationnel. Ce chapitre doit donc être lu comme la mise en place du premier contrat d'usage du système : entrer dans la plateforme, être reconnu, être orienté vers le bon espace et être bloqué dès que l'on sort du périmètre autorisé.

\section{Sprint 1 : objectif et périmètre}

\subsection{Objectif du sprint}
L'objectif principal est de construire la couche d'accès du produit. Dans le contexte de \projectname{}, cela signifie :
\begin{itemize}
  \item authentifier les utilisateurs autorisés par email et mot de passe ;
  \item rediriger automatiquement chaque utilisateur vers le dashboard correspondant à son rôle ;
  \item protéger les routes, layouts et vues métiers ;
  \item fournir un premier niveau d'administration des comptes ;
  \item préparer un cadre de navigation commun pour les modules suivants.
\end{itemize}

Cette séparation entre \textbf{TA}, \textbf{Manager}, \textbf{HR} et \textbf{Admin} n'est pas cosmétique. Elle est réutilisée ensuite dans les pages métiers, dans les Server Actions, dans les services IA et dans l'agent conversationnel, qui doivent tous respecter le périmètre de l'utilisateur connecté.

Au-delà de l'authentification, ce sprint vise aussi trois garanties durables : une confidentialité minimale des données RH, une cohérence de navigation selon le rôle métier et une extensibilité suffisante pour brancher les futurs modules sur un shell commun sans redéfinir la logique d'accès.
Le tableau \ref{tab:access-identity-backlog} détaille les user stories qui traduisent ces objectifs dans le sprint.

\subsection{Backlog du sprint}
\begin{table}[H]
  \centering
  \caption{\label{tab:access-identity-backlog}Backlog fonctionnel du module d'accès et d'identité}
  \begin{tabularx}{\textwidth}{>{\bfseries}p{1.3cm}p{3cm}X>{\raggedright\arraybackslash}p{2cm}}
    \toprule
    ID & Thème & User Story & Statut \\
    \midrule
    1 & Accès sécurisé & En tant qu'utilisateur autorisé, je peux accéder à la plateforme à travers un point d'entrée contrôlé. & Réalisé \\
    2 & Authentification & En tant qu'utilisateur autorisé, je peux me connecter avec mon email et mon mot de passe. & Réalisé \\
    3 & Redirection par rôle & En tant qu'utilisateur connecté, je suis redirigé vers mon espace TA, Manager, HR ou Admin. & Réalisé \\
    4 & Protection des routes & En tant qu'utilisateur non authentifié ou non autorisé, je suis bloqué hors des espaces protégés. & Réalisé \\
    5 & Administration des comptes & En tant qu'administrateur, je peux consulter les utilisateurs, créer un compte, modifier un rôle et bannir ou débannir un utilisateur. & Réalisé \\
    6 & Shell de dashboard & En tant qu'utilisateur connecté, je bénéficie d'un cadre de navigation commun préparant l'accès aux modules métier et au chat analytique. & Réalisé \\
    \bottomrule
  \end{tabularx}
\end{table}

Ce backlog montre déjà que l'identité traverse toute la plateforme. Elle ne se limite pas au formulaire de connexion ; elle englobe la redirection, la protection des routes, la structure des dashboards et la gouvernance des comptes.

\section{Implémentation du sprint}

\subsection{Analyse du sprint}
Le comportement du module repose sur quelques briques fonctionnelles fondamentales :
\begin{itemize}
  \item l'écran de connexion ;
  \item le module d'authentification, de récupération de session et de contrôle d'accès par rôle ;
  \item le layout protégé des espaces dashboard ;
  \item le shell partagé des espaces métier ;
  \item l'interface d'administration des utilisateurs ;
  \item la première vue métier du tableau de bord TA branchée sur ce socle.
\end{itemize}

Ce module ne doit donc pas être lu comme un simple travail de design d'interface. Il formalise le premier contrat d'usage du système : qui peut entrer, où il atterrit, et à quel périmètre il a droit.

D'un point de vue architectural, le sprint combine trois niveaux de contrôle complémentaires : l'identification de l'utilisateur, la vérification de son rôle et la gouvernance des comptes côté administration. Cette superposition réduit le risque qu'une protection repose uniquement sur l'interface.

\subsection{Vue fonctionnelle globale}
Les usages couverts par ce sprint se regroupent en trois blocs :
\begin{itemize}
  \item se connecter ;
  \item être redirigé vers un espace adapté au rôle ;
  \item administrer les comptes si l'utilisateur dispose du rôle administrateur.
\end{itemize}

Le layout de dashboard prépare déjà l'usage futur de l'IA à deux niveaux : une assistance contextuelle disponible dans les espaces métier et un espace dédié aux analyses longues. L'entrée conversationnelle de la plateforme a donc été pensée dès le départ comme une fonctionnalité authentifiée, cloisonnée par rôle et intégrée à la navigation latérale.
Ce point est important : l'IA n'est pas ajoutée plus tard dans un espace parallèle. Elle s'inscrit dès les fondations dans la même logique d'authentification, de layout et de séparation des rôles que le reste du produit.

\subsection{Authentification et redirection par rôle}
L'écran de connexion constitue le premier point de contrôle de la plateforme. Lorsqu'un utilisateur saisit ses identifiants, le système vérifie son identité, récupère son rôle et l'oriente vers l'espace correspondant. Le principe de redirection est volontairement simple : un profil TA rejoint l'espace de recrutement opérationnel, un Manager rejoint son espace d'évaluation, un RH rejoint l'espace de décision finale et un administrateur rejoint l'espace de supervision.

Cette logique n'est pas seulement appliquée dans l'interface. Elle est également vérifiée côté serveur afin que la sécurité ne dépende pas de la navigation du navigateur. Les routes sensibles restent protégées par une lecture réelle de la session et du rôle.
La figure \ref{fig:auth-role-sequence} formalise cette séquence d'authentification et de redirection par rôle.

\begin{figure}[H]
  \centering
  \diagraminput{figures/diagrams/auth-role-sequence}
  \caption{\label{fig:auth-role-sequence}Séquence d'authentification et de redirection par rôle}
\end{figure}

\subsection{Protection du shell et contrôle d'accès}
Lorsqu'un utilisateur accède à un espace protégé, le cadre principal du dashboard vérifie d'abord qu'une session valide existe. En l'absence de session, l'utilisateur est renvoyé vers la connexion. Si la session existe, le rôle détermine les vues accessibles, les actions autorisées et les données visibles.

Cette centralisation apporte trois bénéfices majeurs. D'une part, elle évite de répéter la logique d'accès dans chaque page. D'autre part, elle garantit une expérience homogène pour tous les rôles. Enfin, elle prépare les futurs composants transverses à s'exécuter dans le même contexte de sécurité que les vues métiers.

\subsection{Administration des comptes}
L'administration des comptes complète la couche d'identité. Elle permet à l'organisation de consulter les utilisateurs, de créer les accès nécessaires, de modifier les rôles et de suspendre un compte lorsque le contexte l'exige. Cette gouvernance est essentielle dans une plateforme RH : les personnes autorisées à consulter ou manipuler des candidatures doivent elles-mêmes être gérées de manière explicite.

La possibilité de suspendre ou réactiver un compte apporte également une souplesse importante. Elle permet de retirer rapidement un accès sans supprimer brutalement l'utilisateur ni perturber l'ensemble de la structure des rôles.

\paragraph{Scénarios de validation fonctionnelle.}
La cohérence du module d'accès peut être vérifiée par trois scénarios simples. Premièrement, un utilisateur authentifié doit être redirigé vers l'espace correspondant à son rôle. Deuxièmement, un utilisateur sans session doit être renvoyé vers la page de connexion lorsqu'il tente d'accéder à un dashboard protégé. Troisièmement, les actions d'administration des comptes doivent rester réservées au rôle administrateur. Ces scénarios matérialisent le contrat principal du sprint : aucune donnée métier sensible ne doit être exposée sans session ni rôle approprié.

\section{Réalisation}


\subsection{Page de connexion}
La page de connexion propose un parcours direct : saisie de l'email, saisie du mot de passe, contrôle visuel du champ sensible, état de chargement et message d'erreur compréhensible. Elle confirme aussi que le produit n'est pas pensé comme un portail ouvert au public, mais comme une plateforme d'entreprise dont les comptes sont gérés par l'organisation.

Sur le plan UX, le formulaire reste volontairement sobre. L'objectif est de réduire l'ambiguïté au moment de l'entrée dans la plateforme et de rendre visible le caractère administré de l'accès. La figure \ref{fig:sign-in-page} présente cette interface de connexion.

\begin{figure}[H]
  \centering
  \appcapture{figures/screenshots/sign-in.png}
  \caption{\label{fig:sign-in-page}Page de connexion sécurisée de la plateforme.}
\end{figure}

\subsection{Shell de dashboard partagé}
Le shell de dashboard constitue un apport structurant du sprint. Il mutualise la navigation latérale, l'en-tête, le contexte utilisateur et l'enveloppe commune de rendu des pages métier. Grâce à ce shell, les modules suivants peuvent s'intégrer progressivement sans reconstruire à chaque fois les fondations de navigation et de sécurité.

En pratique, ce shell joue le rôle de charnière entre sécurité et productivité. Il offre un cadre stable pour toutes les vues TA, Manager, RH et Admin, tout en garantissant qu'elles partagent les mêmes conventions d'identité et de navigation.

\subsection{Première vue métier : tableau de bord TA}
La page de tableau de bord TA illustre la première exploitation de ce socle. Elle présente les indicateurs utiles au recruteur et le planning opérationnel du jour, tout en respectant le périmètre du rôle connecté. Ce tableau de bord montre que le sprint ne s'arrête pas à l'identité : il ouvre déjà la porte au pilotage métier réel.

Ce point est important pour la lecture du rapport : l'accès protégé n'est pas une fin en soi. Il débouche immédiatement sur une valeur métier visible, ce qui prouve que le socle construit est déjà exploitable dans l'usage quotidien.
La figure \ref{fig:ta-dashboard} montre cette première surface métier après authentification.

\begin{figure}[H]
  \centering
  \appcapture{figures/screenshots/dashboard-ta.png}
  \caption{\label{fig:ta-dashboard}Tableau de bord TA après authentification.}
\end{figure}

\subsection{Gestion des utilisateurs côté administration}
La page d'administration des utilisateurs matérialise la gouvernance des accès. L'administrateur peut consulter les comptes, créer un utilisateur, modifier un rôle et suspendre ou réactiver un accès selon le besoin.

Cette fonctionnalité matérialise un principe de responsabilité : dans une plateforme RH, les utilisateurs autorisés à voir ou manipuler les candidatures doivent eux-mêmes être gérés de manière explicite et traçable.

\section{Apport du module}
Ce premier module apporte plusieurs fondations durables au projet :
\begin{itemize}
  \item un point d'entrée authentifié et contrôlé pour le produit ;
  \item une authentification alignée sur un usage d'entreprise ;
  \item une séparation explicite par rôle, réutilisée ensuite dans les routes, les services et l'IA ;
  \item un shell commun qui accélère tous les développements suivants ;
  \item un environnement authentifié, contrôlé et traçable pour les futurs usages analytiques.
\end{itemize}

Autrement dit, même si ce module ne contient pas encore le c\oe{}ur du pipeline RAG, il construit l'infrastructure d'usage sans laquelle une IA de recrutement ne pourrait pas être exposée de manière sûre.

On peut donc le considérer comme la fondation de confiance du produit. Sans cette base, les modules suivants auraient pu exister techniquement, mais pas comme un système crédible dans un cadre professionnel exigeant.

\section*{Conclusion}
\addcontentsline{toc}{section}{Conclusion}
Ce premier module établit les fondations d'accès et d'identité de \projectname{}. Il transforme une simple connexion en point d'entrée vers une plateforme gouvernée par les rôles, protégée côté serveur et prête à accueillir les fonctionnalités métier et IA des modules suivants. Le chapitre suivant s'appuie sur ce socle pour traiter le premier c\oe{}ur documentaire et opérationnel du produit : le catalogue de CV et la gestion des offres d'emploi.
